\documentclass[10pt,letterpaper]{article}
\usepackage[margin=0.85in]{geometry}
\usepackage[T1]{fontenc}
\usepackage{lmodern}
\usepackage{amsmath,booktabs}
\usepackage[hidelinks]{hyperref}
\newcommand{\claim}[1]{\textsuperscript{\scriptsize[#1]}}
\title{Development Evidence Update: Source Semantics and Scorer Stability}
\author{Exploratory addendum to the v0.7 methods draft}
\date{2 October 2026 \quad | \quad v0.8}
\begin{document}
\maketitle

\noindent\textbf{Status.} This addendum reports bounded development work toward
an ACL 2027 research paper. It supplements, without rewriting, the v0.7 methods
skeleton. Claim markers map to \texttt{claim\_ledger\_v08.json}. Neither source
annotation nor the scorer control establishes held-out accuracy, an algorithmic
advantage, or novelty.

\section{From discovery records to supplied-document prefixes}
\label{sec:capture}
We attempted all 39 URLs in the frozen 19-history discovery frame. Thirty-two
responses were provisionally usable, covering 15 complete source sets. These
counts describe capture feasibility, not 15 independent evaluation examples;
overlapping organizations and events still require grouping. Four histories
were selected before annotation: IBM, Docker, Chrome, and HS2. Failed IBM
capture blocked both planned prefixes without replacement. The remaining three
histories supplied six cumulative prefixes under a frozen delivery manifest.
Delivery order is author-defined, and retrieval, reported publication, and
event times remain distinct. Current pages do not certify historical original
versions or public availability.\claim{V08-C01}

This distinction materially changes interpretation. Docker's earlier-dated
page already contains a later reversal update, so its first supplied prefix
cannot represent pre-reversal knowledge. The HS2 landing page contains a
correction notice but not a separately captured original erroneous version.
Neither stream permits inventing an earlier information state by deleting
later edits. Captured bytes and extraction hashes support provenance; they do
not resolve these historical limits.\claim{V08-C02}

\section{Source-grounded development annotation}
\label{sec:annotation}
Separate fresh-context model roles produced six source-only annotations and
six reference artifacts. The annotations contain 57 claim instances across
prefixes; repeated cumulative evidence means these are not 57 independent
facts. Two fixed questions per history were reused across its two prefixes,
yielding 12 reference judgments from six unique question texts. Annotation
roles did not receive reference questions, and reference roles did not receive
annotation outputs. All model revisions remain unknown. Procedural separation
does not establish independent training, human gold, or exhaustive candidate
coverage.\claim{V08-C03}

Subsequent source-only development adjudication changed four of the six
annotation artifacts while retaining all reference judgments. Unsupported
Docker public-image correction records were removed; a policy qualification
or clarification does not automatically entail an earlier assertion that must
be withdrawn. Docker reversal timing also required an open lower bound.
Chrome revisions separated an asserted but undated banner from dated article
content and removed an unsupported temporal ordering. The six selected
annotations still contain 57 claim instances. All 12 selected annotation and
reference artifacts pass structural validators; one raw annotation had failed
a date-precision rule. Structural validity remains separate from semantic
entailment.\claim{V08-C04}

HS2 supports an explicitly reported correction from 81\% to 41\%, with a
narrowed compensation-scheme scope. The quoted erroneous value is not a
current-world assertion. Its later route record supports the Manchester plan
change but does not establish the complete Leeds disposition required by one
question. These cases motivate careful treatment of correction targets,
modality, scope, and abstention. They do not yet demonstrate uncertain identity,
cross-document correction propagation, or a need for joint decoding. The
current annotations are source-semantic records, not automatically scored
candidate graphs for the existing decoder.\claim{V08-C05}

\section{Controlling the scorer's execution history}
\label{sec:control}
The preceding v0.7 pilot observed 26 support-score sign changes among 32 items
when answer-definition order changed. Its paired execution reused prompt
prefixes, confounding wording with numerical execution history. The follow-up
selects the first original-order item of each kind: link, unresolved reading,
resolved reading, and null link. This fixed subset was chosen after aggregate
v0.7 instability was known. All four items share one development history; this
is a technical diagnosis, not an independent four-example benchmark.\claim{V08-C06}

For each item and answer order, a fresh process first evaluates the complete
original prompt without cached tokens, then repeats the identical token array
with prefix reuse. Eight processes yield 16 distributions using the same
pinned Qwen3-0.6B Q8\_0 weights and llama.cpp runtime as v0.7. The score remains
$\ell=\log P(\mathrm{Yes})-\log P(\mathrm{No})$ at the same unconstrained
next-token position, using full-vocabulary probabilities. No label-only
normalization, prompt averaging, or score-based item removal is introduced.
Both literal-label probabilities are retained.\claim{V08-C07}

All four cold comparisons change score sign (Table~\ref{tab:control}).
Absolute cold order differences range from 3.211 to 8.132 natural-log units.
All eight immediate warm repeats have exactly matching archived Yes/No log
probabilities, satisfying the fixed $10^{-6}$ tolerance. A separate replay
independently recomputed every measurement and checked the native process
bindings. Prompt-cache carryover is therefore not a sufficient explanation
for these four observed order effects. This does not explain the remaining
28 v0.7 items, prove general numerical invariance, or distinguish broader
prompt sensitivity from poor task calibration.\claim{V08-C08}

\begin{table}[ht]
\centering
\small
\begin{tabular}{lrrrr}
\toprule
Item / kind & Cold YN & Cold NY & Warm--cold YN & Warm--cold NY \\
\midrule
s0001 / Link & +5.601 & -0.517 & 0.000 & 0.000 \\
s0003 / Unresolved & +4.386 & -2.175 & 0.000 & 0.000 \\
s0004 / Reading & +7.010 & -1.121 & 0.000 & 0.000 \\
s0006 / Null link & +3.095 & -0.116 & 0.000 & 0.000 \\
\bottomrule
\end{tabular}
\caption{Observed log odds, rounded to three decimals. YN and NY denote
Yes/No definition order; label meanings remain fixed. Warm--cold columns
subtract the corresponding cold score from its identical-prompt warm repeat.
All rows share one previously inspected history.}
\label{tab:control}
\end{table}


\section{Implications for the next study}
\label{sec:next}
Scorer robustness and source semantics are separate empirical gates. A
correction-looking passage can fail the intended correction relation, while
a structurally valid record can still misstate scope or time. Development
adjudication must therefore precede method comparisons, with unchanged raw
outputs and disclosed repairs. The next evaluation needs grouped source
histories, human or independently adjudicated references, measured candidate
coverage, a stronger pinned contextual baseline, and identical candidates and
scores across decoders. Optimization gaps must be reported separately from
source-grounded question accuracy.\claim{V08-C09}

The distributed record binds code, manifests, results, and review artifacts.
Full captured text and verbatim working annotations have separate
redistribution constraints; hashes alone cannot reconstruct missing texts.
This addendum is LaTeX source only: the available TeX installation lacks the
required format/classes, so no compiled or visually verified PDF is claimed.
The prior methods draft and bibliography remain unchanged.\claim{V08-C10}
\end{document}
